\chapter{DENEYSEL YÖNTEM}

\section{Yazılım/Çalışma Ortamı}
Deneyler Windows üzerinde Python ve ONNX Runtime kullanılarak yürütülmüştür.
Çalışma zamanı, öncelikli olarak OpenVINO EP kullanacak şekilde yapılandırılmış; uyumsuzluk durumlarında CPU EP'ye düşüşe izin verilmiştir.
Sağlayıcı kullanılabilirliği ve cihaz derleme davranışı sistemden sisteme değişebileceği için, pipeline her koşuda kullanılan sağlayıcı zincirini raporlamaktadır.

Deneylerde kullanılan sistem özellikleri aşağıdaki gibidir:
\begin{itemize}
	\item \textbf{Cihaz}: MATEBOOK
	\item \textbf{İşlemci}: Intel(R) Core(TM) Ultra 5 125H (3.60\,GHz)
	\item \textbf{Bellek}: 16\,GB RAM
	\item \textbf{İşletim sistemi}: Windows 11 Home Single Language (25H2), derleme 26200.8037 (64-bit, x64)
\end{itemize}

\section{Modeller}
Üç temsili ONNX modeli seçilmiştir:
\begin{itemize}
	\item \textbf{MobileNetV2}: kompakt CNN.
	\item \textbf{ResNet50}: daha derin CNN.
	\item \textbf{BERT-tiny}: Transformer tabanlı encoder.
\end{itemize}
Bu seçim, farklı operatör karışımlarını ve ölçekleri kapsar.

\section{Ölçüm Metodolojisi}
Her model ve konfigürasyon için tekrarlı koşular gerçekleştirilmiştir.
Her koşuda sabit sayıda çıkarım iterasyonu, ısınma (warm-up) sonrası ölçülmüştür.
Ortalama gecikme (ms), standart sapma ve yaklaşık 95\% güven aralığı raporlanmıştır.

\section{Baseline ve Optimized Konfigürasyonları}
Baseline konfigürasyonda ORT grafik optimizasyonları kapatılmıştır.
Optimized konfigürasyonda extended optimizasyonlar açılmıştır \cite{onnxruntime}.
Optimizasyon sonrası ONNX dışa aktarımının EP bağımsız seri hale getirilebilir kalması için, optimize model dışa aktarımı CPU EP üzerinden yapılmıştır.
Bu tercih, bazı EP'lerin ürettiği derlenmiş düğümler (compiled nodes) nedeniyle seri hale getirme hatalarını engeller.

\section{Üretilen Çıktılar ve Tekrarlanabilirlik}
Pipeline çıktıları aşağıdaki artefaktları üretir:
\begin{itemize}
	\item \texttt{results/performance\_summary.csv}: tek model gecikme özeti.
	\item \texttt{results/scalability\_matrix.csv}: çoklu model matrisi.
	\item \texttt{results/*\_profiling.json}: profiling izleri.
	\item \texttt{paper/figures/}: raporda kullanılan görseller (otomatik kopyalanır).
\end{itemize}